% Part: incompleteness
% Chapter: representability-in-q
% Section: sigma1-completeness

\documentclass[../../../include/open-logic-section]{subfiles}

\begin{document}

\olfileid[hi]{inc}{inp}{s1c}
\olsection{\texorpdfstring{$\Sigma_1$}{Sigma-1} पूर्णता}

$\Th{Q}$ और उसके संगत, अभिगृहीतीकरणीय विस्तार अपूर्ण होने के बावजूद हमने
देखा है कि $\Th{Q}$ अंक-सूचकों के बारे में अनेक मूलभूत तथ्य सिद्ध करता है।
वास्तव में इसे काफी अधिक विस्तृत किया जा सकता है। यह समझने के लिए कि
$\Th{Q}$ में क्या-क्या सिद्ध किया जा सकता है, हम $\Delta_0$, $\Sigma_1$
और $\Pi_1$ !!{formula}s की धारणाएँ प्रस्तुत करते हैं। मोटे तौर पर,
$\Sigma_1$ !!{formula},~$\lexists[x][!B(x)]$ रूप का होता है, जहाँ $!B$ केवल
प्रतिज्ञप्तिक संयोजकों और परिबद्ध परिमाणकों से बनाया गया है। हम दिखाएँगे कि
यदि $!A$ ऐसा $\Sigma_1$ !!{sentence} है जो $\Struct{N}$ में सत्य है, तो
$\Th{Q} \Proves !A$ (\olref{thm:sigma1-completeness})।

\begin{defn}
\ollabel{defn:bd-quant}
\emph{परिबद्ध अस्तित्वात्मक !!{formula}} वह है जो
$\lexists[x][(x < t \land !A(x))]$ के रूप का हो, जहाँ $t$ कोई भी पद है;
इसे हम परंपरानुसार $\bexists{x < t}{!A(x)}$ लिखते हैं।
%
\emph{परिबद्ध सार्विक !!{formula}} वह है जो
$\lforall[x][(x < t \lif !A(x))]$ के रूप का हो, जहाँ $t$ कोई भी पद है;
इसे हम परंपरानुसार $\bforall{x < t}{!A(x)}$ लिखते हैं।
\end{defn}

\begin{defn}
\ollabel{defn:delta0-sigma1-pi1-frm}
कोई !!{formula} $!B$,~$\Delta_0$ है यदि वह परमाण्विक !!{formula}s से केवल
प्रतिज्ञप्तिक संयोजकों और परिबद्ध परिमाणीकरण द्वारा बना हो।
%
कोई !!{formula} $!A$,~$\Sigma_1$ है यदि
$!A \ident \lexists[x][!B(x)]$, जहाँ $!B$,~$\Delta_0$ है।
%
कोई !!{formula} $!A$,~$\Pi_1$ है यदि
$!A \ident \lforall[x][!B(x)]$, जहाँ $!B$,~$\Delta_0$ है।
\end{defn}

\begin{lem}
\ollabel{lem:q-proves-clterm-id} मान लें कि $t$ ऐसा बंद पद है कि
$\Value{t}{N} = n$। तब $\Th{Q} \Proves \eq[t][\num n]$।
\end{lem}

\begin{proof}
हम इसे~$t$ की जटिलता पर आगमन द्वारा सिद्ध करते हैं। आधार मामले में
$\Value{\Obj 0}{N} = 0$, और $\num 0 \ident \Obj 0$ होने से
$\Th{Q} \Proves \eq[\Obj 0][\num 0]$।
%
आगमन मामले के लिए $t_1$ और $t_2$ ऐसे पद हों कि
$\Value{t_1}{N} = n_1$, $\Value{t_2}{N} = n_2$,
$\Th{Q} \Proves \eq[t_1][\num n_1]$ और
$\Th{Q} \Proves \eq[t_2][\num n_2]$।

तब $\Value{(t_1')}{N} = n_1 + 1$, और आगमन परिकल्पना तथा !!{formula}
$\eq[\num{n_1}'][\num{n_1}']$ पर समता के प्रथम-कोटि नियम लागू करने से
$\Th{Q} \Proves \eq[t_1'][{\num n_1}']$; अतः अंक-सूचकों की परिभाषा से
$\Th{Q} \Proves \eq[t_1'][\num{n_1 + 1}]$।

योगों के लिए हमारे पास
\[
      \Value{(t_1 + t_2)}{N}
    = \Value{t_1}{N} + \Value{t_2}{N}
    = n_1 + n_2.
\]
आगमन परिकल्पना और समता के नियमों से
$\Th{Q} \Proves \eq[t_1 + t_2][\num{n_1} + t_2]$, और फिर समता के नियमों
के दूसरे अनुप्रयोग से
$\Th{Q} \Proves \eq[t_1 + t_2][\num{n_1} + \num{n_2}]$।
\olref[inc][req][bre]{lem:q-proves-add} से
$\Th{Q} \Proves \eq[\num{n_1} + \num{n_2}][\num{n_1 + n_2}]$, इसलिए
$\Th{Q} \Proves \eq[t_1 + t_2][\num{n_1 + n_2}]$।

इसी प्रकार का तर्क~$\times$ के लिए भी काम करता है, जिसमें
\olref[inc][req][bre]{lem:q-proves-mult} का उपयोग होता है।
%
इससे अंकगणित के सभी बंद पद समेट लिए जाते हैं, अतः प्रत्येक ऐसे बंद पद~$t$
के लिए जिसके लिए $\Value{t}{N} = n$, हमारे पास
$\Th{Q} \Proves \eq[t][\num n]$ है।
\end{proof}

\begin{prob}
\olref{lem:q-proves-clterm-id} के~$\times$ वाले भाग को विस्तार से सिद्ध
करें।
\end{prob}

\begin{lem}
\ollabel{lem:atomic-completeness}
मान लें कि $t_1$ और $t_2$ बंद पद हैं। तब
\begin{enumerate}
\item यदि $\Value{t_1}{N} = \Value{t_2}{N}$, तो
    $\Th{Q} \Proves \eq[t_1][t_2]$।
\item यदि $\Value{t_1}{N} \neq \Value{t_2}{N}$, तो
    $\Th{Q} \Proves \eq/[t_1][t_2]$।
\item यदि $\Value{t_1}{N} < \Value{t_2}{N}$, तो
    $\Th{Q} \Proves t_1 < t_2$।
\item यदि $\Value{t_2}{N} \leq \Value{t_1}{N}$, तो
    $\Th{Q} \Proves \lnot(t_1 < t_2)$।
\end{enumerate}
\end{lem}

\begin{proof}
पद $t_1$ और $t_2$ दिए जाने पर हम $n = \Value{t_1}{N}$ और
$m = \Value{t_2}{N}$ नियत करते हैं।

मान लें $!A \ident t_1 = t_2$। \olref{lem:q-proves-clterm-id} से
$\Th{Q} \Proves \eq[t_1][\num n]$ और
$\Th{Q} \Proves \eq[t_2][\num n]$। यदि $n = m$, तो
$\Th{Q} \Proves \eq[\num n][\num m]$, और इसलिए समता की संक्रामकता से
$\Th{Q} \Proves \eq[t_1][t_2]$। यदि $n \neq m$, तो
$\Th{Q} \Proves \eq/[\num n][\num m]$, और समता की संक्रामकता से फिर
$\Th{Q} \Proves \eq/[t_1][t_2]$।

अब $!A \ident t_1 < t_2$ मानें। दोनों मामलों के लिए हम अभिगृहीत~$!Q_8$
पर निर्भर करते हैं, जो सभी $x,y$ के लिए
$x < y \liff \lexists[z][\eq[z' + x][y]]$ कहता है।

मान लें $\Sat{N}{t_1 < t_2}$। तब कोई $k \in \Nat$ ऐसा है कि
$n + k + 1 = m$। \olref{lem:q-proves-clterm-id} से
$\Th{Q} \Proves \eq[t_1][\num n]$ और
$\Th{Q} \Proves \eq[t_2][\num m]$, और इस उपपत्ति के पहले भाग से
$\Th{Q} \Proves \eq[\num n + {\num k}'][\num m]$। समता की संक्रामकता
से निकलता है कि $\Th{Q} \Proves \eq[{\num k}' + t_1][t_2]$, इसलिए
$\Th{Q} \Proves \lexists[z][\eq[z' + t_1][t_2]]$। $!Q_8$ की
दाएँ-से-बाएँ दिशा से $\Th{Q} \Proves t_1 < t_2$।

इसके स्थान पर मान लें $\Sat/{N}{t_1 < t_2}$, अर्थात् $m \leq n$।
%
हम~$\Th{Q}$ के भीतर काम करते हुए $t_1 < t_2$ मानते हैं। $!Q_8$ की
बाएँ-से-दाएँ दिशा से कोई~$z$ ऐसा है कि $\eq[z' + t_1][t_2]$। चूँकि
$\Th{Q} \Proves \eq[t_1][\num n]$ और
$\Th{Q} \Proves \eq[t_2][\num m]$, इसलिए
$\eq[z' + \num n][\num m]$।
%
$!Q_5$ का उपयोग करते हुए~$m$ पर बाहरी आगमन से
$\eq[z' + \num{n - m}][\Obj 0]$। यदि $m = n$, तो
$\eq/[z'][\Obj 0]$, जिससे~$!Q_3$ द्वारा विरोधाभास मिलता है। यदि $m < n$,
तो फिर~$!Q_5$ से $\eq[(z' + \num{n - m - 1})'][\Obj 0]$, जिससे~$!Q_3$
द्वारा विरोधाभास मिलता है। अतः $\Th{Q} \Proves \lnot(t_1 < t_2)$।
\end{proof}

\begin{lem}
\ollabel{lem:bounded-quant-equiv}
मान लें कि $!A$ कोई !!a{formula}, $t$ कोई बंद पद और $k=\Value{t}{N}$ है।
तब
\begin{enumerate}
\item $\Th{Q} \Proves \bforall{x<t}{!A(x)}$ तभी और केवल तभी जब
    $\Th{Q} \Proves !A(\num 0) \land \dots \land !A(\num{k-1})$।
\item $\Th{Q} \Proves \bexists{x<t}{!A(x)}$ तभी और केवल तभी जब
    $\Th{Q} \Proves !A(\num 0) \lor \dots \lor !A(\num{k-1})$।
\end{enumerate}
\end{lem}

\begin{proof}
    हम परिबद्ध सार्विक परिमाणक का मामला सिद्ध करते हैं। यदि
    $\Value{t}{N} = 0$, तो समतुल्यता का बायाँ पक्ष~$\Th{Q}$ में सिद्ध है,
    क्योंकि \olref[inc][req][min]{lem:less-zero} के अनुसार कोई
    $x<\num 0$ नहीं है। इसी प्रकार हम रिक्त वियोजन को केवल $\ltrue$ मान
    सकते हैं, जो~$\Th{Q}$ में भी सिद्ध है।
    %
    अतः हम मानते हैं कि किसी प्राकृतिक संख्या~$k$ के लिए
    $\Value{t}{N} = k+1$। \olref{lem:q-proves-clterm-id} से हम मान सकते
    हैं कि हम $\bforall{x<\num{k+1}}{!A(x)}$ रूप के !!a{formula} के साथ
    काम कर रहे हैं।

    मान लें कि $\Th{Q} \Proves \bforall{x<\num{k+1}}{!A(x)}$, और
    $n \leq k$। चूँकि \olref{lem:atomic-completeness} से
    $\Th{Q} \Proves \num n < \num{k+1}$, तर्क से निकलता है कि
    $\Th{Q} \Proves !A(\num n)$। प्रत्येक $n \leq k$ के लिए इस तथ्य को
    $k+1$ बार लागू करके हमें इच्छित
    $\Th{Q} \Proves !A(\num 0) \land \dots \land !A(\num k)$ मिलता है।

    दूसरी दिशा के लिए मान लें कि
    $\Th{Q} \Proves !A(\num 0) \land \dots \land !A(\num k)$।
    $\Th{Q}$ में काम करते हुए मान लें कि $x < \num{k+1}$।
    \olref[inc][req][min]{lem:less-nsucc} से हमारे पास
    $x = \num 0 \lor \dots \lor x = \num k$ है, इसलिए तर्क से~$!A(x)$
    निकलता है, और फलतः सार्विक दावा
    $\bforall{x<\num{k+1}}{!A(x)}$ निकलता है।

    परिबद्ध अस्तित्वात्मक रूप से परिमाणित !!{formula}s की समतुल्यता का
    प्रमाण इसी प्रकार है।
\end{proof}

\begin{prob}
\olref{lem:bounded-quant-equiv} के अस्तित्वात्मक मामले का विस्तृत प्रमाण
दें।
\end{prob}

\begin{lem}
\ollabel{lem:delta0-completeness}
यदि $!A$ ऐसा $\Delta_0$ !!{sentence} है जो $\Struct{N}$ में सत्य है, तो
$\Th{Q} \Proves !A$।
\end{lem}

\begin{proof}
हम इसे !!{formula} की जटिलता पर आगमन द्वारा सिद्ध करते हैं।
%
आधार मामला \olref{lem:atomic-completeness} से मिलता है, इसलिए हम आगमन चरण
पर आते हैं। सरलता के लिए निषेध के मामले को उस !!{formula} की संरचना के
आधार पर उपमामलों में बाँटते हैं जिस पर निषेध लागू होता है।

\begin{enumerate}
\item मान लें $(!A \land !B)$,~$\Struct{N}$ में सत्य है; अतः $!A$ और $!B$
दोनों~$\Struct{N}$ में सत्य हैं। आगमन परिकल्पना से
$\Th{Q} \Proves !A$ और $\Th{Q} \Proves !B$, इसलिए तर्क से
$\Th{Q} \Proves (!A \land !B)$।
%
\item मान लें $\lnot (!A \land !B)$,~$\Struct{N}$ में सत्य है; अतः या तो
$\lnot !A$ अथवा $\lnot !B$,~$\Struct{N}$ में सत्य है। सामान्यता की हानि
के बिना पहले को मान लें। आगमन परिकल्पना से
$\Th{Q} \Proves \lnot !A$, और इसलिए तर्क से
$\Th{Q} \Proves \lnot (!A \land !B)$।
%
\item मान लें $(!A \lor !B)$,~$\Struct{N}$ में सत्य है; अतः या तो $!A$,
~$\Struct{N}$ में सत्य है अथवा $!B$,~$\Struct{N}$ में सत्य है। सामान्यता की हानि के बिना पहले को
मान लें। आगमन परिकल्पना से $\Th{Q} \Proves !A$, और इसलिए तर्क से
$\Th{Q} \Proves (!A \lor !B)$।
%
\item मान लें $\lnot(!A \lor !B)$,~$\Struct{N}$ में सत्य है; अतः
$\lnot !A$ और $\lnot !B$,~$\Struct{N}$ में सत्य हैं। तब आगमन परिकल्पना से
$\Th{Q} \Proves \lnot !A$ और $\Th{Q} \Proves \lnot !B$। फलतः तर्क से
$\Th{Q} \Proves \lnot(!A \lor !B)$।
%
\item मान लें $\bforall{x<t}{!A(x)}$,~$\Struct{N}$ में सत्य है, जहाँ $t$
कोई बंद पद और $k=\Value{t}{N}$ है। आगमन परिकल्पना और तर्क से, यदि प्रत्येक
$n < \Value{t}{N}$ के लिए $!A(\num n)$,~$\Struct{N}$ में सत्य है, तो
$\Th{Q} \Proves !A(\num 0) \land \dots \land !A(\num{k-1})$।
\olref{lem:bounded-quant-equiv} से निकलता है कि
$\Th{Q} \Proves \bforall{x<t}{!A(x)}$।
%
\item परिबद्ध अस्तित्वात्मक परिमाणक का मामला, जिसमें हमारे पास
$\bexists{x < t}{!A(x)}$ रूप का !!a{sentence} है, परिबद्ध सार्विक परिमाणक
के मामले जैसा है।
%
\item मान लें $\lnot \bforall{x<t}{!A(x)}$,~$\Struct{N}$ में सत्य है,
जहाँ $t$ कोई बंद पद है। यह !!{sentence},~!!{sentence} $\bexists{x<t}{\lnot !A(x)}$ के
समतुल्य है और यह समतुल्यता~$\Th{Q}$ में व्युत्पन्न की जा सकती है; अतः हम
परिबद्ध अस्तित्वात्मक परिमाणकों वाला तर्क लागू कर सकते हैं।
%
\item इसी प्रकार मान लें कि $\lnot \bexists{x<t}!A(x)$,~$\Struct{N}$ में
सत्य है, जहाँ $t$ कोई बंद पद है। यह !!{sentence},~$\Th{Q}$ में
$\bforall{x<t}{\lnot!A(x)}$ के समतुल्य है, इसलिए हम परिबद्ध सार्विक
परिमाणकों वाला तर्क लागू कर सकते हैं।
%
\item अंततः मान लें $\lnot !A$,~$\Struct{N}$ में सत्य है। अब केवल वे
मामले शेष हैं जिनमें $!A$ परमाण्विक है, अथवा किसी $\Delta_0$
!!{sentence} $!B$ के लिए $\lnot !A \ident \lnot\lnot !B$। यदि $!A$
परमाण्विक है, तो \olref{lem:atomic-completeness} से
$\Th{Q} \Proves \lnot !A$। यदि
$\lnot !A \ident \lnot\lnot !B$, तो तर्क द्वारा यह~$\Th{Q}$ में~$!B$ के
सिद्धतः समतुल्य है; और चूँकि $\lnot !A$,~$\Struct{N}$ में सत्य है, इसलिए
वह भी~$\Struct{N}$ में सत्य है। अतः आगमन परिकल्पना से हमारे पास
$\Th{Q} \Proves \lnot !A$ है।
\end{enumerate}
\end{proof}

\begin{prob}
\olref{lem:delta0-completeness} के अस्तित्वात्मक मामले का विस्तृत प्रमाण
दें।
\end{prob}

\begin{thm}
\ollabel{thm:sigma1-completeness}
यदि $!A$ ऐसा $\Sigma_1$ !!{sentence} है जो~$\Struct{N}$ में सत्य है, तो
$\Th{Q} \Proves !A$।
\end{thm}

\begin{proof}
यदि $\lexists{x}!A(x)$ ऐसा $\Sigma_1$ !!{sentence} है जो~$\Struct{N}$ में
सत्य है, तो कोई प्राकृतिक संख्या~$n$ और कोई चर-नियतन~$s$ ऐसे हैं कि
$s(x) = n$ और $\Sat{N}{!A(x)}[s]$। संतुष्टि-संबंध के मानक तथ्यों से निकलता
है कि $\Sat{N}{!A(\num n)}$। किंतु $!A(\num n)$ कोई $\Delta_0$
!!{formula} है, इसलिए \olref{lem:delta0-completeness} से हमारे पास
$\Th{Q} \Proves !A(\num n)$ है, और फलतः तर्क से हमारे पास
$\Th{Q} \Proves \lexists[x][!A(x)]$ भी है।
\end{proof}

\end{document}
